\documentclass[11pt,a4paper]{article}
\usepackage[margin=2.5cm]{geometry}
\usepackage{amsmath,amssymb,amsthm}
\usepackage{booktabs}
\usepackage{array}
\usepackage{xcolor}
\usepackage{tikz}
\usepackage{pgfplots}
\pgfplotsset{compat=1.18}
\usepackage{hyperref}
\usepackage{parskip}
\usepackage{enumitem}
\usepackage{fancyhdr}
\usepackage{caption}

\definecolor{mipblue}{RGB}{30,90,160}
\definecolor{mipgreen}{RGB}{20,130,70}
\definecolor{mipred}{RGB}{180,30,30}
\definecolor{mipgray}{RGB}{80,80,80}
\definecolor{lightblue}{RGB}{220,235,250}
\definecolor{lightgreen}{RGB}{220,245,225}

\hypersetup{
  colorlinks=true,
  linkcolor=mipblue,
  urlcolor=mipblue,
  citecolor=mipblue,
  pdftitle={MICLSP: Multi-Item Capacitated Lot Sizing with Batch Production},
  pdfauthor={MIPster Test Suite},
}

\pagestyle{fancy}
\fancyhf{}
\fancyhead[L]{\small\color{mipgray}MIPster Test Suite}
\fancyhead[R]{\small\color{mipgray}MICLSP Problem Description}
\fancyfoot[C]{\small\color{mipgray}\thepage}
\renewcommand{\headrulewidth}{0.4pt}

\newcommand{\var}[1]{\mathit{#1}}
\newcommand{\param}[1]{\mathit{#1}}

\begin{document}

%% ────────────────────────────────────────────────────────────────────────────
%%  Title block
%% ────────────────────────────────────────────────────────────────────────────
\begin{center}
  {\LARGE\bfseries\color{mipblue}
    Multi-Item Capacitated Lot Sizing\\[4pt]
    with Batch Production (MICLSP)}\\[0.8em]
  {\large\color{mipgray} MIPster Test Suite --- Problem Description \& Formulation}\\[0.4em]
  {\small\color{mipgray} All three variable types $\cdot$ Big-M constraints $\cdot$
   Parameterisable difficulty}
\end{center}

\vspace{0.5em}
\hrule height 1pt
\vspace{1.5em}

%% ────────────────────────────────────────────────────────────────────────────
%%  1. Motivation
%% ────────────────────────────────────────────────────────────────────────────
\section{Motivation and Problem Class}

The \textbf{Multi-Item Capacitated Lot Sizing Problem with Batch Production}
(MICLSP) is a classical production-planning problem that arises whenever a
shared machine must produce several items in discrete batches over a finite
planning horizon.

It is an excellent benchmark for a mixed-integer programming solver because
it \emph{naturally} contains all three variable types that challenge a B\&B
solver, and the Big-M linking constraint appears organically from the problem
structure rather than as a modelling artefact:

\begin{itemize}[leftmargin=1.8em]
  \item \textbf{Binary variables} $y_{it}$: whether a costly machine setup is
        performed for item~$i$ in period~$t$.
  \item \textbf{General-integer variables} $n_{it}$: the number of production
        batches of item~$i$ in period~$t$ (each batch contains a fixed number
        of units~$B_i$).
  \item \textbf{Continuous variables} $s_{it}$: the inventory of item~$i$ held
        at the end of period~$t$ (fractional storage units arise from rounding
        and cost interpolation).
  \item \textbf{Big-M linking constraint}: a batch can only be produced if the
        machine was set up; the tightest valid coefficient
        $M_{it}=\lfloor C_t/B_i\rfloor$ comes directly from the machine
        capacity in period~$t$.
\end{itemize}

The LP relaxation typically yields a fractional mix of setup and batch
variables, so instances are rarely solved at the root node, making MICLSP
a reliable stress test for the branching and cutting machinery of a MIP solver.

%% ────────────────────────────────────────────────────────────────────────────
%%  2. Data
%% ────────────────────────────────────────────────────────────────────────────
\section{Problem Data}

Let $\mathcal{I}=\{0,\ldots,I{-}1\}$ be the set of items and
$\mathcal{T}=\{0,\ldots,T{-}1\}$ the set of time periods.

\begin{center}
\begin{tabular}{@{}lll@{}}
\toprule
\textbf{Symbol} & \textbf{Domain} & \textbf{Meaning} \\
\midrule
$B_i$   & $\mathbb{Z}_{>0}$ & Batch size of item $i$ (units per batch) \\
$d_{it}$ & $\mathbb{Z}_{\ge 0}$ & Demand for item $i$ in period $t$ \\
$f_i$   & $\mathbb{R}_{>0}$ & Fixed setup cost for item $i$ \\
$h_i$   & $\mathbb{R}_{>0}$ & Holding cost per unit of item $i$ per period \\
$C_t$   & $\mathbb{Z}_{>0}$ & Machine capacity (units) in period $t$ \\
$M_{it}$& $\mathbb{Z}_{\ge 0}$ & Big-M coefficient: $M_{it}=\lfloor C_t/B_i\rfloor$ \\
\bottomrule
\end{tabular}
\end{center}

%% ────────────────────────────────────────────────────────────────────────────
%%  3. Decision variables
%% ────────────────────────────────────────────────────────────────────────────
\section{Decision Variables}

\begin{center}
\begin{tabular}{@{}lll@{}}
\toprule
\textbf{Variable} & \textbf{Type} & \textbf{Meaning} \\
\midrule
$y_{it}$  & Binary ($\{0,1\}$) & 1 if item $i$ is set up in period $t$ \\
$n_{it}$  & Integer ($\mathbb{Z}_{\ge 0}$) & Number of batches of item $i$ produced in period $t$ \\
$s_{it}$  & Continuous ($\mathbb{R}_{\ge 0}$) & Inventory of item $i$ at the end of period $t$ \\
\bottomrule
\end{tabular}
\end{center}

The convention $s_{i,-1}=0$ (no initial stock) is used throughout.

%% ────────────────────────────────────────────────────────────────────────────
%%  4. Mathematical formulation
%% ────────────────────────────────────────────────────────────────────────────
\section{Mathematical Formulation}

\begin{equation}\label{eq:obj}
\min \quad \sum_{i\in\mathcal{I}}\sum_{t\in\mathcal{T}}
  \Bigl( f_i\, y_{it} \;+\; h_i\, s_{it} \Bigr)
\end{equation}

\noindent subject to:

\begin{alignat}{3}
  &s_{i,t-1} + B_i\, n_{it} - s_{it} = d_{it}
  &&\quad\forall\, i\in\mathcal{I},\; t\in\mathcal{T}
  &&\quad\text{(inventory balance)}\label{eq:inv}\\[4pt]
  &\sum_{i\in\mathcal{I}} B_i\, n_{it} \le C_t
  &&\quad\forall\, t\in\mathcal{T}
  &&\quad\text{(machine capacity)}\label{eq:cap}\\[4pt]
  &n_{it} \le M_{it}\, y_{it}
  &&\quad\forall\, i\in\mathcal{I},\; t\in\mathcal{T}
  &&\quad\text{(Big-M setup link)}\label{eq:bigm}\\[4pt]
  &y_{it} \in \{0,1\}
  &&\quad\forall\, i\in\mathcal{I},\; t\in\mathcal{T}
  &&\\[2pt]
  &n_{it} \in \mathbb{Z}_{\ge 0}
  &&\quad\forall\, i\in\mathcal{I},\; t\in\mathcal{T}
  &&\\[2pt]
  &s_{it} \ge 0
  &&\quad\forall\, i\in\mathcal{I},\; t\in\mathcal{T}
  &&
\end{alignat}

\paragraph{Sizes.}
An instance with $I$ items and $T$ periods has
$I\cdot T$ binary variables, $I\cdot T$ general-integer variables,
and $I\cdot T$ continuous variables,
for a total of $3\cdot I\cdot T$ variables and
$I\cdot T + T + I\cdot T = T(2I+1)$ constraints
(excluding variable bounds).

%% ────────────────────────────────────────────────────────────────────────────
%%  5. Big-M discussion
%% ────────────────────────────────────────────────────────────────────────────
\section{The Big-M Linking Constraint}

Constraint~\eqref{eq:bigm} is the classic \emph{Big-M} formulation for
linking a semi-continuous variable to a binary indicator.  Here,
$M_{it}=\lfloor C_t/B_i\rfloor$ is the \textbf{tightest possible} coefficient:
it equals the maximum number of batches of item~$i$ that could physically be
produced in period~$t$ given the machine capacity~$C_t$.  Using this tight
value (rather than an arbitrary large constant) strengthens the LP relaxation
and reduces the integrality gap.

If $y_{it}=0$, then $n_{it}\le 0$, so $n_{it}=0$ (no production without
setup).  If $y_{it}=1$, the constraint $n_{it}\le M_{it}$ is redundant,
since the capacity constraint~\eqref{eq:cap} already limits the total batches
across all items.

\colorbox{lightblue}{\parbox{0.96\textwidth}{%
\textbf{Why Big-M is interesting here.}
The LP relaxation of~\eqref{eq:bigm} sets $n_{it}\le M_{it}\,\hat{y}_{it}$
where $\hat{y}_{it}\in[0,1]$.  A fractional setup $\hat{y}_{it}=0.4$ allows
fractional batch counts up to $0.4\,M_{it}$, which combined with
fractional inventory $\hat{s}_{it}$ gives the LP solver large freedom to
reduce setup costs while satisfying demands.  This fractional exploitation is
precisely what the B\&B solver must overcome by branching on $y_{it}$.
}}

%% ────────────────────────────────────────────────────────────────────────────
%%  6. LP relaxation and integrality gap
%% ────────────────────────────────────────────────────────────────────────────
\section{LP Relaxation and Integrality Gap}

The LP relaxation is obtained by replacing
$y_{it}\in\{0,1\}$ with $0\le y_{it}\le 1$ and
$n_{it}\in\mathbb{Z}_{\ge 0}$ with $n_{it}\ge 0$.
Let $z^*_{\mathrm{LP}}$ and $z^*_{\mathrm{MIP}}$ denote the optimal values.
Since this is a minimisation problem, $z^*_{\mathrm{LP}}\le z^*_{\mathrm{MIP}}$.

The \textbf{relative integrality gap} is defined as
\[
  \Delta\% \;=\; 100\;\cdot\;
  \frac{z^*_{\mathrm{MIP}} - z^*_{\mathrm{LP}}}{\max(1,\,z^*_{\mathrm{LP}})}.
\]

Instances selected as test fixtures satisfy $\Delta\% > 0$, confirming that
the root-node LP bound is strictly weaker than the MIP optimum and that
B\&B (with cuts and heuristics) is required to close the gap.

%% ────────────────────────────────────────────────────────────────────────────
%%  7. Random instance generation
%% ────────────────────────────────────────────────────────────────────────────
\section{Random Instance Generation}

Instances are generated deterministically from a seed using a linear
congruential generator with the following distributions:

\begin{center}
\begin{tabular}{@{}llll@{}}
\toprule
\textbf{Parameter} & \textbf{Distribution} & \textbf{Range} & \textbf{Interpretation}\\
\midrule
$B_i$ & Uniform integer & $[5, 30]$ & Batch size \\
$d_{it}$ & Uniform integer & $[1, 20]$ & Period demand \\
$f_i$ & Uniform integer & $[100, 500]$ & Setup cost \\
$h_i$ & Uniform integer & $[1, 10]$ & Holding cost \\
$C_t$ & Derived & $\lceil\rho\cdot\sum_i d_{it}\rceil$ & Machine capacity \\
\bottomrule
\end{tabular}
\end{center}

The \textbf{capacity tightness ratio}
$\rho = C_t \,/\, \sum_i d_{it}$
is the key difficulty knob.  With $\rho>1$ the machine can accommodate all
current-period demand on average, but the discrete batch structure and setup
sequencing force careful scheduling.  Typical values used in the test suite:

\begin{center}
\begin{tabular}{@{}cll@{}}
\toprule
$\rho$ & \textbf{Effect} & \textbf{Typical solve} \\
\midrule
$\ge 2.0$ & Very loose; most instances solved at root & $< 1\,\text{s}$, $\Delta\%\approx 0$ \\
$1.5$     & Moderate; some branching required & $1$--$10\,\text{s}$\\
$1.25$--$1.35$ & Tight; clear LP-MIP gap, hundreds of nodes & $10\,\text{s}$--$2\,\text{min}$ \\
$1.10$--$1.20$ & Very tight; thousands of nodes & $1$--$5\,\text{min}$ \\
$< 1.0$   & Often infeasible & --- \\
\bottomrule
\end{tabular}
\end{center}

\noindent
For the test suite we target the $\rho=1.25$--$1.35$ band: instances
solve in $10\,\text{s}$--$2\,\text{min}$ with a non-trivial LP-MIP gap,
exercising branching, cuts (MIR, GMI), and FeasibilityJump.

%% ────────────────────────────────────────────────────────────────────────────
%%  8. Example: small instance
%% ────────────────────────────────────────────────────────────────────────────
\section{Illustrative Example ($I=2$, $T=3$)}

Consider two items ($I=2$) and three periods ($T=3$) with the following data:

\begin{center}
\begin{tabular}{@{}ccccc@{}}
\toprule
Item $i$ & $B_i$ & $f_i$ & $h_i$ & $d_{it}$ for $t=0,1,2$ \\
\midrule
0 & 10 & 200 & 2 & 8, 5, 12 \\
1 & 15 & 150 & 3 & 6, 9,  4 \\
\bottomrule
\end{tabular}
\qquad
\begin{tabular}{@{}cc@{}}
\toprule
$t$ & $C_t$ \\
\midrule
0 & 30 \\
1 & 30 \\
2 & 30 \\
\bottomrule
\end{tabular}
\end{center}

Variables: $y_{it}\in\{0,1\}$, $n_{it}\in\mathbb{Z}_{\ge 0}$,
$s_{it}\in\mathbb{R}_{\ge 0}$.  Tight Big-M values:
$M_{i0}=\lfloor 30/B_i\rfloor$, i.e.\
$M_{00}=3$, $M_{10}=2$ (similarly for other periods).

\smallskip
\textbf{Optimal solution.}  In $t=0$: set up both items
($y_{00}=y_{10}=1$), produce 2 batches of item~0 (20 units) and 1 batch of
item~1 (15 units), using capacity $10{\cdot}2+15{\cdot}1=35>30$.
A feasible schedule might produce item~0 in $t=0,2$ and item~1 in $t=0,1$:

\begin{center}
\begin{tabular}{@{}c|ccc|ccc|ccc@{}}
\toprule
 & \multicolumn{3}{c|}{$t=0$} & \multicolumn{3}{c|}{$t=1$} & \multicolumn{3}{c}{$t=2$}\\
 & $y$ & $n$ & $s$ & $y$ & $n$ & $s$ & $y$ & $n$ & $s$\\
\midrule
Item 0 & 1 & 2 & 12 & 0 & 0 & 7 & 1 & 1 & 5\\
Item 1 & 1 & 1 & 9  & 1 & 1 & 1 & 0 & 0 & 0\\% example only; not verified optimal
\bottomrule
\end{tabular}
\end{center}

This example illustrates how the three variable types interact: the binary
$y$ triggers a setup cost, the integer $n$ determines batch production, and
the continuous $s$ accumulates the remaining inventory.

%% ────────────────────────────────────────────────────────────────────────────
%%  9. Solution verification
%% ────────────────────────────────────────────────────────────────────────────
\section{Solution Verification in the Test Suite}

After solving, the test binary verifies the following for every saved
incumbent solution:

\begin{enumerate}[leftmargin=2em]
  \item \textbf{Inventory balance} (equality constraint): for each $(i,t)$,
    $\lvert s_{i,t-1} + B_i\,n_{it} - s_{it} - d_{it}\rvert \le \varepsilon$.
  \item \textbf{Machine capacity} (inequality): $\sum_i B_i\,n_{it}\le C_t+\varepsilon$.
  \item \textbf{Big-M / no-production-without-setup}:
    if $n_{it}>0$ then $y_{it}=1$.
  \item \textbf{Non-negativity}: $s_{it}\ge -\varepsilon$, $n_{it}\ge -\varepsilon$.
  \item \textbf{Integrality}: $n_{it}$ is within $\varepsilon$ of an integer;
    $y_{it}\in\{0,1\}$ within $\varepsilon$.
  \item \textbf{Objective match}: the reported objective matches
    $\sum_{i,t}(f_i\,y_{it}+h_i\,s_{it})$ within $\varepsilon$.
  \item \textbf{Known optimal}: the best solution objective matches the
    pre-computed optimal value stored in \texttt{miclsp\_fixtures.h}.
\end{enumerate}

Here $\varepsilon = 10^{-4}$ (the solver primal feasibility tolerance).

%% ────────────────────────────────────────────────────────────────────────────
%%  10. Instance table (filled in by fixture generator)
%% ────────────────────────────────────────────────────────────────────────────
\section{Test Fixture Instances}

The table below lists all instances stored in \texttt{miclsp\_fixtures.h}.
Column ``gap'' is the LP--MIP integrality gap $\Delta\%$; ``B\&B'' indicates
whether at least one branch was performed (gap~$>0$).

\begin{center}
\begin{tabular}{@{}lcccccrl@{}}
\toprule
\textbf{Instance tag} & $I$ & $T$ & $\rho$ & \textbf{Vars} &
  \textbf{Opt} & \textbf{LP bound} & \textbf{Gap}\\
\midrule
\emph{(generated by gen\_miclsp\_fixtures.py)} & & & & & & &\\
\bottomrule
\end{tabular}
\end{center}

The fixture header is regenerated by running:
\begin{verbatim}
  cd Cbc/test
  python3 gen_miclsp_fixtures.py > miclsp_fixtures.h
\end{verbatim}

%% ────────────────────────────────────────────────────────────────────────────
%%  References
%% ────────────────────────────────────────────────────────────────────────────
\section*{References}

\begin{enumerate}[leftmargin=2em,label={[\arabic*]}]
  \item W.~W.~Trigeiro, L.~J.~Thomas, J.~O.~McClain,
    ``Capacitated lot sizing with setup times,''
    \textit{Management Science} 35(3):353--366, 1989.
  \item Y.~Pochet, L.~A.~Wolsey,
    \textit{Production Planning by Mixed Integer Programming},
    Springer, 2006.
  \item C.~P.~M.~van~Hoesel, A.~P.~M.~Wagelmans,
    ``An $O(T^3)$ algorithm for the economic lot-sizing problem
    with constant capacities,''
    \textit{Management Science} 42(1):142--150, 1996.
\end{enumerate}

\end{document}
